\section{Introduction}
\label{ch:introduction}

This blueprint organizes a Lean formalization of the theorem $\MIPstar = \RE$ of
Ji, Natarajan, Vidick, Wright and Yuen~\cite{JNVWY20}. The main result is that
there is a computable map from Turing machines to nonlocal games that sends halting
machines to games of (quantum) value $1$ and non-halting machines to games of value at
most $\frac12$; in particular, approximating the quantum value of a nonlocal game is an
uncomputable problem. Through the work of Fritz~\cite{Fri12}, Junge et
al.~\cite{JNPPSW11} and Ozawa~\cite{Oza13}, this refutes Tsirelson's problem and Connes'
Embedding Problem.

The main theorem is stated in the self-contained (other than importing Mathlib) file
\texttt{HaltingGameValue.lean}. Every definition it contains
(POVMs, synchronous games, game descriptions) is elementary linear algebra and
computability. The goal of this projects are to (a) complete a formalization of this theorem and
(b) do it in such a way that the required effort benefits the community down-the-line, e.g.\ by
making the proof modular, stating and proving as many facts as possible in a manner that is reusable, etc.
Finally, (c) we expect multiple simplifications and clarifications to the proof to accompany the effort.


This initial blueprint is a coarse sketch. It maps out the main steps of the proof, in the
attempt to modularize it and provide a diversity of entrypoints to the formalization effort.
The most important definitions are complete, but most theorems appear only in ``headline'' form,
 with the
decomposition into formalizable lemmas left as a project task.

% The one place
%where the blueprint is intended to be fully precise is the statement of the main theorem
%(Section~\ref{ch:main-theorem}), which matches the self-contained Lean file
%\texttt{HaltingGameValue.lean} stated on top of Mathlib alone.

\subsection{Overview of the proof}
\label{sec:overview}

The inclusion $\MIPstar \subseteq \RE$ is elementary: the value of a game can be
approximated from below by enumerating tensor-product strategies
(Section~\ref{sec:rr-value-approx}). The substance is $\RE \subseteq \MIPstar$, which is proved by
a reduction from the halting problem.
The central ingredient in the proof is a \emph{gap-preserving compression}
procedure for ``normal form'' verifiers (Section~\ref{sec:ps-compression}): a polynomial-time
transformation that maps a verifier $\verifier$ to a verifier $\vcompr$ which, at index
$n$, simulates the exponentially larger game $\verifier_{2^n}$, preserving the value-$1$
property and forcing exponential entanglement lower bounds. Given compression, an
abstract fixed-point argument (the \emph{recursive compression lemma},
Section~\ref{sec:rr-compression}, applied via Kleene's recursion theorem) produces the
halting reduction.

Compression itself is assembled from three transformations, each described in
Section~\ref{ch:proof-structure}:
\emph{question reduction} by introspection (the players sample their own questions,
certified by a Pauli-basis self-test built on the quantum low individual degree test),
\emph{answer reduction} by PCP composition (the players prove, via a probabilistically
checkable proof, that their answers would have been accepted), and \emph{gap
amplification} by anchored parallel repetition. Each of these steps has one central
ingredient that can be formalized in an independent manner: the low individual degree
test, the PCP theorem (in a specific form), and a parallel repetition theorem (for a certain
class of quantum games).

\subsection{Entry points}
\label{sec:entry-points}

Contributors with different scientific backgrounds can enter the project in different places, largely independently.

\emph{Quantum information.} The foundations chapter
(Section~\ref{ch:foundations}) and the required results on self-testing: the Magic
Square game (Section~\ref{sec:rr-magic-square}), the Gowers--Hatami theorem
(Section~\ref{sec:rr-gowers-hatami}), the orthonormalization lemma
(Section~\ref{sec:rr-orthonormalization}), and the low individual degree test
(Section~\ref{sec:rr-qld}). These are self-contained statements about states and
measurements, formalizable without any complexity theory.

\emph{Complexity and computability theory.} The computability toolkit
(Section~\ref{sec:rr-computability}), the abstract recursive compression lemma
(Section~\ref{sec:rr-compression}), succinct descriptions and the Cook--Levin machinery
feeding answer reduction (Section~\ref{sec:rr-cook-levin}). These are classical
statements with no quantum content, and several partially exist in Mathlib.

\emph{Operator algebras.} The synchronous framework of
Section~\ref{ch:foundations}, the bridge between synchronous and tensor-product values
(Section~\ref{sec:rr-almost-sync}), and the downstream consequences
(Section~\ref{ch:downstream}), which connect the main theorem to Tsirelson's problem
and Connes' Embedding Problem.

In addition, contributors can participate at different levels. Contributors with a high degree of Lean expertise may want
to focus on the foundations and concentrate on those elements of the proof that are of broader interest and may later
be upstreamed to Mathlib, Physlib, CSLib, etc. For example, the Gowers Hatami theorem, general facts about Turing machines and computability, the PCP theorem,
formalization of quantum games and quantum game values, etc.

\subsection{Departures from the paper}
\label{sec:departures}

The formalization is not intended to follow~\cite{JNVWY20} line by line. Since the initial publication of the paper
a number of simplifications have appeared in the literature. In particular, we plan to incorporate the following.

\paragraph{Synchronous strategies throughout.}
The paper works with bipartite tensor-product strategies $(\ket{\psi}, A, B)$ on
$\mH_\alice \ot \mH_\bob$. We instead take \emph{synchronous} strategies --- a single
family of measurements $\{M^x_a\}$ on one space, evaluated against the tracial state ---
as the basic object, and define the synchronous value $\synval$
(Definition~\ref{def:sync-value}). This eliminates the duplication of Hilbert spaces,
states, and consistency bookkeeping that pervades the paper's soundness analyses, and it
matches the main Lean statement. The cost is a single bridging step: the
almost-synchronous correlations theorem (Section~\ref{sec:rr-almost-sync}) shows that
for synchronous games the two values agree closely enough to transport completeness and
soundness. All games constructed in the proof are synchronous (or are made so by
symmetrization), so this is a genuine simplification.

\paragraph{Projective strategies via orthonormalization.}
Rather than invoking Naimark dilation to replace POVMs by projective measurements (which
enlarges the space and interacts poorly with the tracial framework), we use the
orthonormalization lemma (Section~\ref{sec:rr-orthonormalization}) to replace
almost-projective POVMs by nearby projective measurements on the \emph{same} space. All
soundness analyses should be carried out for projective strategies from the start.

\paragraph{Computable rather than polynomial-time, where possible.}
The final undecidability statement only needs the halting reduction to be
\emph{computable}; polynomial-time bounds matter only inside the compression recursion,
where the runtime of $\Compress$ must be smaller than the succinctness it exploits (see
the discussion in~\cite{MNY}). The blueprint keeps time bounds where they are load-bearing
(normal form verifiers, compression) and drops them everywhere else. In particular the
main theorem (Section~\ref{ch:main-theorem}) is stated with a computable reduction.

\paragraph{Abstract recursive compression.}
The fixed-point argument that turns compression into a halting reduction is stated once,
abstractly, following Marks--Nezhadi--Yuen~\cite{MNY}
(Section~\ref{sec:rr-compression}), and instantiated with $f = $ entanglement
requirement. This isolates all uses of Kleene's recursion theorem in one
quantum-free lemma.

\paragraph{Other candidate simplifications (to be decided).}
Stating soundness of all transformations directly in terms of the entanglement measure
$\Ent$ rather than via separate value and dimension statements; using de la Salle's
quantitative improvements to stability (Gowers--Hatami) where they shorten proofs; other
simplifications that are present in the follow-up proof of
$\MIPco=\coRE$~\cite{Lin23,Lin25}.

\subsection{Conventions}
\label{sec:conventions}

Every definition and theorem carries a \verb|\label|, and \verb|\uses| annotations
record the dependency graph rendered on the project web page. Lean declaration names
(\verb|\lean| tags) are deliberately absent at this stage; they will be added as the
corresponding declarations are created. Statements in
Sections~\ref{ch:required-results}--\ref{ch:anchored-repetition} are external results: the project may initially
axiomatize them (each section records the literature source and comments on
formalization), while statements in Section~\ref{ch:proof-structure} are proof
obligations of the project. Numerical constants in headline statements (levels, error
exponents) are copied from~\cite{JNVWY20} and should be treated as provisional until the
corresponding section is worked out in detail.
